\chapter{I processi software}

\section{Cos'è un processo software}

Nel capitolo precedente abbiamo visto che serve un approccio \textbf{sistematico} alla produzione del software.
Il primo strumento per ottenerlo è il concetto di \textit{processo}.

\dfn{Processo software}{
  Un \textbf{processo software} è un insieme di \textbf{attività correlate} che porta alla produzione di un sistema software.
}

Un'analogia utile è la costruzione di una casa: non si comincia posando mattoni a caso, ma si segue una sequenza
di attività (capire le esigenze di chi ci abiterà, fare il progetto, ottenere i permessi, costruire, collaudare, fare manutenzione).
Il processo software è l'equivalente per il software: stabilisce \textbf{quali attività} svolgere, \textbf{in che ordine},
\textbf{chi} le svolge e \textbf{cosa} producono.

\warning{Non esiste il processo perfetto}{
  Non esiste un processo universale che funzioni sempre. Esistono (e si usano) \textbf{molti processi diversi}, e ognuno va
  \textbf{adattato alla situazione}: al tipo di software, all'azienda, al cliente, al team.
}

Per capire perché, basta confrontare due progetti molto diversi:

\begin{table}[H]
\centering
\renewcommand{\arraystretch}{1.4}
\begin{tabularx}{\textwidth}{>{\bfseries}P{4cm} L L}
\toprule
 & \textcolor{swred}{Software di un pacemaker} & \textcolor{swblue}{App di una startup} \\
\midrule
Conseguenze di un errore & Possono essere letali & Utenti insoddisfatti, perdita di clienti \\
Stabilità dei requisiti & Molto stabili, fissati da normative & Cambiano di settimana in settimana \\
Documentazione & Estesa e obbligatoria (certificazioni) & Ridotta al minimo indispensabile \\
Priorità & Sicurezza e correttezza & Velocità di rilascio e feedback \\
Processo adatto & Rigoroso, pianificato, molto controllato & Leggero, iterativo, flessibile \\
\bottomrule
\end{tabularx}
\caption{Contesti diversi richiedono processi diversi.}
\end{table}

% ============================================================
\section{Le quattro attività fondamentali}

Per quanto i processi possano essere diversi tra loro, \textbf{tutti} includono, in qualche forma, quattro attività fondamentali.

\begin{figure}[H]
\centering
\begin{tikzpicture}
  \foreach \i/\titolo/\col/\desc in {
      0/{Software\\Specification}/swteal/{Si definiscono le \textbf{funzionalità} e i \textbf{vincoli operativi} del software.},
      1/{Software\\Implementation}/swblue/{Si \textbf{produce} il software che soddisfa i requisiti.},
      2/{Software\\Validation}/ph3/{Si \textbf{verifica} che il software soddisfi i requisiti.},
      3/{Software\\Evolution}/ph6/{Il software \textbf{evolve} per rispondere a esigenze che cambiano.}}{
    \node[phase=\col, minimum width=3.4cm, minimum height=1.3cm, font=\small\bfseries, drop shadow] (a\i) at (\i*4.55,0) {\titolo};
    \node[note, text width=3.5cm, align=center, font=\small, anchor=north] at (\i*4.55,-0.9) {\desc};
  }
  \foreach \i [evaluate=\i as \j using int(\i+1)] in {0,1,2}{
    \node[single arrow, fill=swgray!40, draw=none, minimum height=0.9cm, minimum width=0.7cm, single arrow head extend=1.5mm, inner sep=0pt] at ($(a\i.east)!0.5!(a\j.west)$) {};
  }
\end{tikzpicture}
\caption{Le quattro attività fondamentali presenti in ogni processo software.}
\end{figure}

\subsection{Software Specification (specifica)}

È l'attività in cui si stabilisce \textbf{cosa} il sistema deve fare e \textbf{sotto quali vincoli} deve operare.
Si parla con il cliente e con gli utenti, si analizza il dominio del problema e si mettono per iscritto i requisiti.

\begin{esempio}{Specifica di un sistema per la prenotazione degli esami}
  \textit{Funzionalità}: lo studente può vedere gli appelli disponibili, prenotarsi e cancellare la prenotazione.\\
  \textit{Vincoli operativi}: il sistema deve gestire 5.000 prenotazioni contemporanee nel giorno di apertura degli appelli
  e deve essere accessibile anche da smartphone.
\end{esempio}

\subsection{Software Implementation (sviluppo)}

È l'attività in cui si \textbf{produce} il software che soddisfa la specifica. Attenzione: nonostante il nome, non si tratta
solo di scrivere codice. L'implementazione in senso ampio comprende la \textbf{progettazione} (come organizzare il sistema)
e la \textbf{programmazione} vera e propria. In alcuni testi questa attività è infatti chiamata \textit{Software Development}.

\subsection{Software Validation (validazione)}

Il software prodotto va \textbf{controllato}: bisogna assicurarsi che faccia quello che la specifica richiede
e che risponda alle aspettative del cliente. Questo avviene principalmente tramite revisioni e \textbf{test}.

\subsection{Software Evolution (evoluzione)}

Il software \textbf{non è mai «finito»}. Una volta in uso deve cambiare per rispondere a esigenze che cambiano. Alcuni esempi tipici:
\begin{itemize}
  \item esce una \textbf{nuova versione del sistema operativo} e l'applicazione va adattata;
  \item viene scoperto un \textbf{bug} sfuggito ai test;
  \item entra in vigore una \textbf{nuova legge} che cambia le procedure (es.\ un nuovo calcolo delle tasse);
  \item il cliente chiede \textbf{nuove funzionalità}.
\end{itemize}

\info{{Attività complesse, fatte da persone}}{
  \begin{itemize}
    \item Queste attività sono \textbf{complesse} e ciascuna comprende molte \textbf{sotto-attività} (le vedremo nel modello a cascata).
    \item Come ogni processo creativo, i processi software si basano su \textbf{persone} che prendono \textbf{decisioni} ed esprimono \textbf{giudizi}:
          nessun processo può essere eseguito «meccanicamente».
    \item Oltre alle quattro attività principali, un processo include tipicamente \textbf{attività di supporto}, come la
          \textbf{pianificazione} del progetto e il suo \textbf{monitoraggio} (stiamo rispettando tempi e costi?).
  \end{itemize}
}

\begin{figure}[H]
\centering
\begin{tikzpicture}
  \foreach \i/\titolo/\col in {0/Specifica/swteal, 1/Implementazione/swblue, 2/Validazione/ph3, 3/Evoluzione/ph6}{
    \node[phase=\col, minimum width=3.6cm, minimum height=0.9cm, font=\small\bfseries] (b\i) at (\i*4.1,0) {\titolo};
  }
  \node[fit=(b0)(b3), draw=swgray, dashed, rounded corners=4pt, inner sep=6pt, label={[font=\scriptsize\bfseries, text=swgray]above:Attività fondamentali}] (fund) {};
  \node[rectangle, rounded corners=3pt, fill=swgray!15, draw=swgray!60, minimum width=16.1cm, minimum height=0.8cm, font=\small, below=0.35cm of fund] (supp) {\textbf{Attività di supporto}: pianificazione del progetto, monitoraggio di tempi e costi, \ldots};
  \node[font=\small\itshape, text=swgray, below=0.15cm of supp] {\iuser\ \iuser\ \iuser\quad svolte da persone che prendono decisioni e giudizi\quad\iuser\ \iuser\ \iuser};
\end{tikzpicture}
\caption{Le attività di supporto accompagnano per tutta la durata le attività fondamentali.}
\end{figure}

% ============================================================
\section{Il ciclo di vita del software}

Le quattro attività fondamentali, ripetute nel tempo, descrivono il \hi{ciclo di vita} del software.
Si parla di «ciclo» perché l'evoluzione riporta continuamente a nuove specifiche: ogni modifica richiede di capire
cosa cambiare, implementarlo, validarlo e rilasciarlo.

\begin{figure}[H]
\centering
\begin{tikzpicture}
  \def\R{2.6}
  \node[circle, fill=primary!10, draw=primary!60, minimum size=2.9cm, align=center, font=\small\bfseries, text=primary!80!black] (c) at (0,0) {Ciclo di vita\\del software};
  \node[phase=swteal, minimum width=3.1cm] (s) at (90:\R) {Specifica};
  \node[phase=swblue, minimum width=3.1cm] (i) at (0:\R+1.6) {Implementazione};
  \node[phase=ph3, minimum width=3.1cm] (v) at (270:\R) {Validazione};
  \node[phase=ph6, minimum width=3.1cm] (e) at (180:\R+1.6) {Evoluzione};
  \draw[flow, line width=1.5pt, swgray!80] (s.east) to[out=0, in=90] (i.north);
  \draw[flow, line width=1.5pt, swgray!80] (i.south) to[out=270, in=0] (v.east);
  \draw[flow, line width=1.5pt, swgray!80] (v.west) to[out=180, in=270] (e.south);
  \draw[flow, line width=1.5pt, swred] (e.north) to[out=90, in=180] node[above left, font=\scriptsize, text=swred, align=right]{nuove esigenze,\\nuovi requisiti} (s.west);

  % Nascita e morte
  \node[card=swgreen, font=\scriptsize, text width=3.2cm, align=left] (nasc) at (7.6,2.1) {\textbf{Nascita}: qualcuno \textbf{richiede} il software (un cliente o il mercato).};
  \node[card=swred, font=\scriptsize, text width=3.2cm, align=left] (mor) at (7.6,-2.1) {\textbf{Morte}: esplicita attività di \textbf{dismissione}, ad esempio quando viene sostituito da un altro software.};
  \draw[-{Stealth}, swgreen, line width=1pt, dashed] (nasc.west) to[out=180, in=20] (s.east);
  \draw[-{Stealth}, swred, line width=1pt, dashed] (e.south west) to[out=240, in=180] ([yshift=-1.0cm]v.south) -- ([yshift=-1.0cm]v.south -| mor.west) -- ([yshift=-0.2cm]mor.west);
\end{tikzpicture}
\caption{Il ciclo di vita: dalla richiesta iniziale fino alla dismissione.}
\end{figure}

\dfn{Nascita e morte del software}{
  La vita di un software \textbf{inizia} quando qualcuno ne fa richiesta e \textbf{termina} solo con un'esplicita
  \textbf{attività di dismissione} (ad esempio quando va sostituito con un altro sistema). Tra questi due estremi il software
  attraversa ripetutamente le attività di specifica, implementazione, validazione ed evoluzione.
}

Un aspetto spesso sottovalutato è \textbf{quanto} dura questa vita. Lo sviluppo iniziale può durare mesi,
ma l'esercizio e la manutenzione possono durare \textbf{anni o decenni}: molte banche e pubbliche amministrazioni usano
ancora sistemi (detti \textit{legacy}) scritti in COBOL decenni fa.

\begin{figure}[H]
\centering
\begin{tikzpicture}[x=1cm]
  \draw[line width=1pt, -{Stealth}, swgray] (-0.2,0) -- (17.6,0) node[right, font=\scriptsize]{tempo};
  \fill[swteal!70] (0,0.15) rectangle (1.2,0.85);
  \fill[swblue!70] (1.2,0.15) rectangle (4.2,0.85);
  \fill[ph3!70] (4.2,0.15) rectangle (5.2,0.85);
  \fill[ph6!65] (5.2,0.15) rectangle (16.4,0.85);
  \node[font=\tiny\bfseries, text=white] at (0.6,0.5) {Spec.};
  \node[font=\scriptsize\bfseries, text=white] at (2.7,0.5) {Sviluppo};
  \node[font=\tiny\bfseries, text=white] at (4.7,0.5) {Test};
  \node[font=\scriptsize\bfseries, text=white] at (10.8,0.5) {Esercizio e manutenzione (anni, a volte decenni)};
  \foreach \x/\lab/\col in {0/{richiesta}/swgreen, 5.2/{rilascio}/swblue, 16.4/{dismissione}/swred}{
    \draw[\col, line width=1.2pt] (\x,-0.15) -- (\x,1.05);
    \node[font=\scriptsize\bfseries, text=\col, below] at (\x,-0.15) {\lab};
  }
  \draw[decorate, decoration={brace, amplitude=5pt, mirror}, swgray] (0.05,-0.75) -- node[below=5pt, font=\scriptsize]{mesi} (5.15,-0.75);
\end{tikzpicture}
\caption{Durata qualitativa delle fasi della vita di un software: la manutenzione è di gran lunga la più lunga.}
\end{figure}

% ============================================================
\section{Quanto costa il software}

\subsection{Hardware contro software}

Agli albori dell'informatica la parte costosa di un sistema era l'\textbf{hardware}. Con il passare dei decenni il costo
dell'hardware è crollato, mentre il software è diventato sempre più grande e complesso. Oggi la situazione è ribaltata:
la quota maggiore del costo di un sistema è il \textbf{software}.

\begin{figure}[H]
\centering
\begin{tikzpicture}
\begin{axis}[
  width=11.5cm, height=6.6cm,
  xmin=1970, xmax=2000, ymin=0, ymax=100,
  xtick={1970,1980,1990,2000},
  x tick label style={/pgf/number format/1000 sep=},
  ytick={0,25,50,75,100},
  yticklabel={\pgfmathprintnumber{\tick}\%},
  xlabel={Anno}, ylabel={Quota del costo totale},
  label style={font=\small}, tick label style={font=\scriptsize},
  axis on top, axis line style={swgray},
]
  \fill[swgray!18] (axis cs:1970,0) rectangle (axis cs:2000,100);
  \addplot[domain=1970:2000, samples=80, fill=swblue!35, draw=swblue, line width=1.4pt]
     {10+80/(1+exp(-(x-1987)/3.8))} \closedcycle;
  \node[font=\bfseries, text=swgray!90!black] at (axis cs:1977,72) {Hardware};
  \node[font=\bfseries, text=swblue!80!black] at (axis cs:1994,38) {Software};
\end{axis}
\end{tikzpicture}
\caption{Andamento qualitativo dei costi relativi di hardware e software (da D. Bell, \textit{Software Engineering for Students}).}
\end{figure}

\subsection{Dove vanno i soldi: i costi per fase}

Ancora più interessante è vedere \textbf{come} si distribuiscono i costi tra le varie fasi della vita di un software.

\begin{figure}[H]
\centering
\begin{tikzpicture}
  \def\tmW{17}\def\tmH{7}
  \pgfmathsetmacro{\tmM}{\tmW*0.67}
  \pgfmathsetmacro{\tmD}{\tmW-\tmM}
  \pgfmathsetmacro{\tmT}{\tmH*15/32}
  \pgfmathsetmacro{\tmB}{\tmH-\tmT}
  \pgfmathsetmacro{\tmA}{\tmD*6/17}
  \pgfmathsetmacro{\tmBB}{\tmD*6/17}
  % Manutenzione
  \fill[swteal] (0,0) rectangle (\tmM,\tmH);
  \node[anchor=north west, text=white, font=\large\bfseries] at (0.15,\tmH-0.1) {Manutenzione};
  \node[anchor=south west, text=white, font=\Huge\bfseries] at (0.2,0.2) {67\%};
  % Sviluppo
  \fill[sworange] (\tmM,\tmB) rectangle (\tmW,\tmH);
  \node[anchor=north west, text=white, font=\small\bfseries] at (\tmM+0.12,\tmH-0.1) {Sviluppo (33\%)};
  \node[text=white, font=\small\bfseries, align=center] at (\tmM+\tmD/2,\tmB+\tmT/2-0.2) {Testing\\\Large 15\%};
  \fill[sworange!85!black] (\tmM,0) rectangle (\tmM+\tmA,\tmB);
  \fill[sworange!70!black] (\tmM+\tmA,0) rectangle (\tmM+\tmA+\tmBB,\tmB);
  \fill[sworange!55!black] (\tmM+\tmA+\tmBB,0) rectangle (\tmW,\tmB);
  \node[text=white, font=\scriptsize\bfseries, align=center] at (\tmM+\tmA/2,\tmB/2) {Requisiti\\\large 6\%};
  \node[text=white, font=\scriptsize\bfseries, align=center] at (\tmM+\tmA+\tmBB/2,\tmB/2) {Coding\\\large 6\%};
  \node[text=white, font=\scriptsize\bfseries, align=center] at ({\tmM+\tmA+\tmBB+(\tmW-\tmM-\tmA-\tmBB)/2},\tmB/2) {Design\\\large 5\%};
  \draw[white, line width=1.5pt] (\tmM,0) -- (\tmM,\tmH) (\tmM,\tmB) -- (\tmW,\tmB) (\tmM+\tmA,0) -- (\tmM+\tmA,\tmB) (\tmM+\tmA+\tmBB,0) -- (\tmM+\tmA+\tmBB,\tmB);
\end{tikzpicture}
\caption{Costi relativi delle fasi di un processo software: l'area di ogni rettangolo è proporzionale al costo (da D. Bell).}
\end{figure}

Questo grafico contiene alcune delle lezioni più importanti dell'intero corso:

\begin{enumerate}
  \item \textbf{La manutenzione costa circa il doppio di tutto lo sviluppo messo insieme.} Il grosso della spesa arriva
        \textit{dopo} il rilascio.
  \item \textbf{Scrivere codice è una frazione piccola del totale} (circa il 6\%). Chi pensa che «fare software» significhi
        «programmare» sta guardando solo una piccola parte del problema.
  \item \textbf{Il testing costa più della programmazione}: verificare che il software funzioni è un lavoro enorme.
\end{enumerate}

\info{La conseguenza pratica: progettare per il cambiamento}{
  Se due terzi del costo sono manutenzione, allora conviene investire nelle fasi iniziali (requisiti e progettazione)
  per ottenere un software \textbf{facile da modificare}: modulare, leggibile, ben documentato e ben testato.
  Ogni ora risparmiata «scrivendo codice di fretta» si paga con gli interessi negli anni di manutenzione.
}

\warning{Prendere i numeri con le pinze}{
  Le percentuali esatte variano molto da progetto a progetto e da studio a studio. Quello che conta è l'\textbf{ordine di grandezza}:
  la manutenzione domina, la scrittura del codice è una parte minoritaria.
}

% ============================================================
\section{I modelli di processo}

Un processo software reale è pieno di dettagli: persone, documenti, riunioni, strumenti, eccezioni. Per ragionarci sopra
servono delle rappresentazioni semplificate.

\dfn{Modello di processo software (SDLC)}{
  Un \textbf{Software Process Model}, detto anche \textbf{Software Development Life Cycle (SDLC)}, è una
  \textbf{rappresentazione semplificata} di un processo software. Un modello può concentrarsi su una particolare
  \textbf{prospettiva} (ad esempio la sequenza delle attività, i ruoli coinvolti o i documenti prodotti).
}

Il rapporto tra modello e processo è lo stesso che c'è tra una \textbf{cartina} e il \textbf{territorio}: la cartina
non riporta ogni albero, ma è proprio questa semplificazione a renderla utile. E come esistono cartine diverse per scopi diversi
(stradale, topografica, dei sentieri), esistono modelli di processo diversi che mettono in luce aspetti diversi.

Nel corso vedremo alcuni modelli di processo molto generali:

\begin{figure}[H]
\centering
\begin{tikzpicture}
  % Plan-driven
  \node[chip=swblue, minimum width=7.8cm, minimum height=0.7cm, font=\small\bfseries] (t1) at (0,0) {Modelli guidati dal piano (\textit{plan-driven})};
  \foreach \i/\c in {0/ph1,1/ph2,2/ph3,3/ph4,4/ph5,5/ph6}{
    \fill[\c] (-3.6+\i*1.25,-1.0-\i*0.32) rectangle ++(1.05,-0.42);
  }
  \foreach \i in {0,...,4}{
    \draw[-{Stealth[length=1.5mm]}, swgray] (-2.55+\i*1.25,-1.21-\i*0.32) -| (-2.35+\i*1.25,-1.32-\i*0.32);
  }
  \node[note, text width=7.6cm, align=center, font=\small] at (0,-3.95) {Tutto viene pianificato all'inizio; le fasi si susseguono in ordine.\\\textbf{Esempio}: il modello a cascata (\textit{Waterfall}) --- \textit{prossimo capitolo}.};

  % Incremental / agile
  \node[chip=sworange, minimum width=7.8cm, minimum height=0.7cm, font=\small\bfseries, right=1.2cm of t1] (t2) {Modelli incrementali e agili};
  \foreach \k in {0,1,2}{
    \begin{scope}[shift={(6.3+\k*2.6,-1.75)}]
      \draw[sworange, line width=1.3pt, -{Stealth[length=2mm]}] (0:0.75) arc (0:320:0.75);
      \node[font=\scriptsize\bfseries, text=sworange!80!black] at (0,0) {v\the\numexpr\k+1\relax};
    \end{scope}
  }
  \draw[-{Stealth}, swgray, line width=1pt] (7.1,-1.75) -- (7.8,-1.75);
  \draw[-{Stealth}, swgray, line width=1pt] (9.7,-1.75) -- (10.4,-1.75);
  \node[note, text width=7.6cm, align=center, font=\small] at (8.9,-3.95) {Il software cresce per piccoli incrementi, ripetendo le attività in cicli brevi.\\\textbf{Esempio}: modelli incrementali, metodologie agili --- \textit{più avanti nel corso}.};
\end{tikzpicture}
\caption{Due grandi famiglie di modelli di processo.}
\end{figure}

\gbox{In sintesi}{azure-gradient-6!80!black}{
\begin{itemize}
  \item Un \textbf{processo software} è un insieme di attività correlate che porta alla produzione di software. Non ne esiste uno universale.
  \item Ogni processo comprende \textbf{specifica, implementazione, validazione ed evoluzione}, più attività di supporto come pianificazione e monitoraggio.
  \item Queste attività formano il \textbf{ciclo di vita}: dalla richiesta alla dismissione, con la manutenzione come fase più lunga.
  \item La \textbf{manutenzione} è la voce di costo principale (circa 2/3): bisogna progettare pensando al cambiamento.
  \item Un \textbf{modello di processo (SDLC)} è una rappresentazione semplificata di un processo. Il primo che studiamo è il modello a cascata.
\end{itemize}
}
